HESEL-modellen er en sammensætning af fire koblede PDE'er. De beskriver, partikel tætheden \ref{tæthedsligning}, vorticiteten \ref{vorticiteten}, og elektron \ref{elektrontryksligning} og iontrykket \ref{iontryksligning} i kanten af et fusionplasma.

De har formen:
\begin{align}
    d_t n + n,\mathcal{K}(\phi) - \mathcal{K}(p_e) = \Lambda_n + \Sigma_n \label{tæthedsligning} \\
    \nabla \cdot d_t^0 \nabla_\perp \phi^* - \mathcal{K}(p_e+p_i)=\Lambda_\omega+\Sigma_\omega \label{vorticiteten} \\
    \frac{3}{2}d_t p_e + \frac{5}{2}p_e \mathcal{K}(\phi) -\frac{5}{2}\mathcal{K}\left(\frac{p_e^2}{n}\right)= \Lambda_{p_e}+\Sigma_{p_e} \label{elektrontryksligning} \\
    \frac{3}{2}d_t p_i + \frac{5}{2}p_i \mathcal{K}(\phi) +\frac{5}{2}\mathcal{K}\left(\frac{p_i^2}{n}\right) * p_i\mathcal{K}(p_e+p_i) = \Lambda_{p_i}+\Sigma_{p_i} \label{iontryksligning}
\end{align}
Operatorene er givet ved:
\begin{align}
    d_t f = \frac{\partial f}{\partial t} + \frac{1}{B}  \left\{ \phi , f \right\}\notag \\
    \left\{ f, g\right\} = \partial_x f ,\partial_y g - \partial_y f ,\partial_x g \notag \\
    \mathcal{K}(f)= -\kappa \frac{\partial f}{\partial y} \notag
\end{align}

Appalations moddeller af HESEL-modellen kommer til at sætte $\Sigma$ termerne til 0 og gøre systemet til isoterm ved at sætte $T_e$ og $T_i$ til konstanter.

Simuleringen og implimenteringen af HESEL-modellen er udført med BOUT++ og BOUT-HESEL hvilket er udviklet af DTU PPFE.

Data behandling og visualisering er udført i python ved hjælp af xbout, numpy og matplotlib.

HESEL-modellen, træning af SciML blvier udført i Python ved hjælp af pytorch.

